\chapter{Unit 6: Real-World Applications (OpenCV, Pygame, Web Scraping, Toolboxes)}

\section{Section A: Multiple Choice Practice Questions (60 Curated Problems)}

\noindent This section contains 60 rigorous multiple-choice questions extracted from university course examination assessments. Review each question, test edge-case syntax, and verify code execution mechanics.
\vspace{3mm}

\noindent\textbf{Q1.} \hfill {\small\color{LPUOrange}\textbf{[OpenCV-based image operations | Easy]}}\\[2pt]
Which Python library is primarily used for computer vision and image processing tasks?

\begin{enumerate}[label=(\alph*),topsep=2pt,itemsep=1pt,partopsep=0pt,parsep=0pt]
    \item \relax Matplotlib
    \item \relax PyGame
    \item \relax OpenCV
    \item \relax BeautifulSoup
\end{enumerate}
\vspace{2mm}

\noindent\textbf{Q2.} \hfill {\small\color{LPUOrange}\textbf{[OpenCV-based image operations | Easy]}}\\[2pt]
Which function is used to read an image from a file in OpenCV?

\begin{enumerate}[label=(\alph*),topsep=2pt,itemsep=1pt,partopsep=0pt,parsep=0pt]
    \item \relax \texttt{cv2.showImage()}
    \item \relax \texttt{cv2.loadFile()}
    \item \relax \texttt{cv2.openImage()}
    \item \relax \texttt{cv2.imread()}
\end{enumerate}
\vspace{2mm}

\noindent\textbf{Q3.} \hfill {\small\color{LPUOrange}\textbf{[OpenCV-based image operations | Easy]}}\\[2pt]
In OpenCV, which color format are images loaded in by default?

\begin{enumerate}[label=(\alph*),topsep=2pt,itemsep=1pt,partopsep=0pt,parsep=0pt]
    \item \relax HSV
    \item \relax BGR
    \item \relax RGB
    \item \relax CMYK
\end{enumerate}
\vspace{2mm}

\noindent\textbf{Q4.} \hfill {\small\color{LPUOrange}\textbf{[OpenCV-based image operations | Easy]}}\\[2pt]
Which OpenCV function displays an image in a window?

\begin{enumerate}[label=(\alph*),topsep=2pt,itemsep=1pt,partopsep=0pt,parsep=0pt]
    \item \relax \texttt{cv2.view()}
    \item \relax \texttt{cv2.render()}
    \item \relax \texttt{cv2.imshow()}
    \item \relax \texttt{cv2.display()}
\end{enumerate}
\vspace{2mm}

\noindent\textbf{Q5.} \hfill {\small\color{LPUOrange}\textbf{[PyGame-based gaming creation | Easy]}}\\[2pt]
PyGame is a Python library mainly used for creating what type of applications?

\begin{enumerate}[label=(\alph*),topsep=2pt,itemsep=1pt,partopsep=0pt,parsep=0pt]
    \item \relax Web servers
    \item \relax Machine learning models
    \item \relax Database systems
    \item \relax Games and multimedia
\end{enumerate}
\vspace{2mm}

\noindent\textbf{Q6.} \hfill {\small\color{LPUOrange}\textbf{[PyGame-based gaming creation | Easy]}}\\[2pt]
Which function must be called to start using PyGame's modules?

\begin{enumerate}[label=(\alph*),topsep=2pt,itemsep=1pt,partopsep=0pt,parsep=0pt]
    \item \relax \texttt{pygame.begin()}
    \item \relax \texttt{pygame.run()}
    \item \relax \texttt{pygame.start()}
    \item \relax \texttt{pygame.init()}
\end{enumerate}
\vspace{2mm}

\noindent\textbf{Q7.} \hfill {\small\color{LPUOrange}\textbf{[PyGame-based gaming creation | Easy]}}\\[2pt]
In PyGame, which object represents the main window or drawing surface?

\begin{enumerate}[label=(\alph*),topsep=2pt,itemsep=1pt,partopsep=0pt,parsep=0pt]
    \item \relax Frame
    \item \relax Panel
    \item \relax Surface
    \item \relax Canvas
\end{enumerate}
\vspace{2mm}

\noindent\textbf{Q8.} \hfill {\small\color{LPUOrange}\textbf{[PyGame-based gaming creation | Easy]}}\\[2pt]
Which PyGame concept is typically used to continuously update the game and handle events?

\begin{enumerate}[label=(\alph*),topsep=2pt,itemsep=1pt,partopsep=0pt,parsep=0pt]
    \item \relax Data pipeline
    \item \relax Game loop
    \item \relax Event stack
    \item \relax Render tree
\end{enumerate}
\vspace{2mm}

\noindent\textbf{Q9.} \hfill {\small\color{LPUOrange}\textbf{[web-scraping-based url access | Easy]}}\\[2pt]
Which Python library is commonly used to send HTTP requests to access a URL?

\begin{enumerate}[label=(\alph*),topsep=2pt,itemsep=1pt,partopsep=0pt,parsep=0pt]
    \item \relax \texttt{requests}
    \item \relax \texttt{turtle}
    \item \relax \texttt{numpy}
    \item \relax \texttt{pygame}
\end{enumerate}
\vspace{2mm}

\noindent\textbf{Q10.} \hfill {\small\color{LPUOrange}\textbf{[web-scraping-based url access | Easy]}}\\[2pt]
Which library is widely used to parse HTML content during web scraping?

\begin{enumerate}[label=(\alph*),topsep=2pt,itemsep=1pt,partopsep=0pt,parsep=0pt]
    \item \relax Seaborn
    \item \relax BeautifulSoup
    \item \relax PyGame
    \item \relax OpenCV
\end{enumerate}
\vspace{2mm}

\noindent\textbf{Q11.} \hfill {\small\color{LPUOrange}\textbf{[web-scraping-based url access | Easy]}}\\[2pt]
What does the term "web scraping" mean?

\begin{enumerate}[label=(\alph*),topsep=2pt,itemsep=1pt,partopsep=0pt,parsep=0pt]
    \item \relax Hosting websites on a server
    \item \relax Extracting data from websites automatically
    \item \relax Designing website layouts
    \item \relax Encrypting web traffic
\end{enumerate}
\vspace{2mm}

\noindent\textbf{Q12.} \hfill {\small\color{LPUOrange}\textbf{[web-scraping-based url access | Easy]}}\\[2pt]
An HTTP status code of \texttt{200} typically indicates what?

\begin{enumerate}[label=(\alph*),topsep=2pt,itemsep=1pt,partopsep=0pt,parsep=0pt]
    \item \relax Access is forbidden
    \item \relax The request succeeded
    \item \relax A server error occurred
    \item \relax The page was not found
\end{enumerate}
\vspace{2mm}

\noindent\textbf{Q13.} \hfill {\small\color{LPUOrange}\textbf{[data visualization | Easy]}}\\[2pt]
Which Python library is one of the most widely used for creating basic plots and charts?

\begin{enumerate}[label=(\alph*),topsep=2pt,itemsep=1pt,partopsep=0pt,parsep=0pt]
    \item \relax Matplotlib
    \item \relax Requests
    \item \relax PyGame
    \item \relax OpenCV
\end{enumerate}
\vspace{2mm}

\noindent\textbf{Q14.} \hfill {\small\color{LPUOrange}\textbf{[data visualization | Easy]}}\\[2pt]
Which chart type is best suited for showing trends in data over time?

\begin{enumerate}[label=(\alph*),topsep=2pt,itemsep=1pt,partopsep=0pt,parsep=0pt]
    \item \relax Line chart
    \item \relax Pie chart
    \item \relax Word cloud
    \item \relax Scatter matrix
\end{enumerate}
\vspace{2mm}

\noindent\textbf{Q15.} \hfill {\small\color{LPUOrange}\textbf{[data visualization | Easy]}}\\[2pt]
Which library is built on top of Matplotlib and provides attractive statistical graphics?

\begin{enumerate}[label=(\alph*),topsep=2pt,itemsep=1pt,partopsep=0pt,parsep=0pt]
    \item \relax BeautifulSoup
    \item \relax Seaborn
    \item \relax Flask
    \item \relax PyGame
\end{enumerate}
\vspace{2mm}

\noindent\textbf{Q16.} \hfill {\small\color{LPUOrange}\textbf{[data visualization | Easy]}}\\[2pt]
Which chart is most appropriate for showing the proportion of parts to a whole?

\begin{enumerate}[label=(\alph*),topsep=2pt,itemsep=1pt,partopsep=0pt,parsep=0pt]
    \item \relax Histogram
    \item \relax Line chart
    \item \relax Pie chart
    \item \relax Box plot
\end{enumerate}
\vspace{2mm}

\noindent\textbf{Q17.} \hfill {\small\color{LPUOrange}\textbf{[discussion on advanced toolboxes in python | Easy]}}\\[2pt]
Which library is commonly used for numerical computing and working with arrays in Python?

\begin{enumerate}[label=(\alph*),topsep=2pt,itemsep=1pt,partopsep=0pt,parsep=0pt]
    \item \relax Requests
    \item \relax PyGame
    \item \relax NumPy
    \item \relax Pillow
\end{enumerate}
\vspace{2mm}

\noindent\textbf{Q18.} \hfill {\small\color{LPUOrange}\textbf{[discussion on advanced toolboxes in python | Easy]}}\\[2pt]
Which Python library is primarily used for data analysis and handling tabular data?

\begin{enumerate}[label=(\alph*),topsep=2pt,itemsep=1pt,partopsep=0pt,parsep=0pt]
    \item \relax OpenCV
    \item \relax Turtle
    \item \relax Pandas
    \item \relax PyGame
\end{enumerate}
\vspace{2mm}

\noindent\textbf{Q19.} \hfill {\small\color{LPUOrange}\textbf{[demonstration of real-world applications | Easy]}}\\[2pt]
Which of the following is a real-world application where Python is commonly used?

\begin{enumerate}[label=(\alph*),topsep=2pt,itemsep=1pt,partopsep=0pt,parsep=0pt]
    \item \relax Manufacturing car engines
    \item \relax Data analysis and machine learning
    \item \relax Painting walls
    \item \relax Baking bread
\end{enumerate}
\vspace{2mm}

\noindent\textbf{Q20.} \hfill {\small\color{LPUOrange}\textbf{[storytelling | Easy]}}\\[2pt]
In the context of data, what does "data storytelling" mainly involve?

\begin{enumerate}[label=(\alph*),topsep=2pt,itemsep=1pt,partopsep=0pt,parsep=0pt]
    \item \relax Encrypting datasets securely
    \item \relax Writing fictional novels with code
    \item \relax Deleting outdated records
    \item \relax Communicating insights using data and visuals
\end{enumerate}
\vspace{2mm}

\noindent\textbf{Q21.} \hfill {\small\color{LPUOrange}\textbf{[OpenCV-based image operations | Medium]}}\\[2pt]
In OpenCV, when you read an image using \texttt{cv2.imread()}, in what channel order are the pixel values stored by default?

\begin{enumerate}[label=(\alph*),topsep=2pt,itemsep=1pt,partopsep=0pt,parsep=0pt]
    \item \relax HSV
    \item \relax BGR
    \item \relax Grayscale
    \item \relax RGB
\end{enumerate}
\vspace{2mm}

\noindent\textbf{Q22.} \hfill {\small\color{LPUOrange}\textbf{[OpenCV-based image operations | Medium]}}\\[2pt]
You want to convert a color image \texttt{img} to grayscale in OpenCV. Which line accomplishes this?

\begin{enumerate}[label=(\alph*),topsep=2pt,itemsep=1pt,partopsep=0pt,parsep=0pt]
    \item \relax \texttt{gray = cv2.imread(img, 0)}
    \item \relax \texttt{gray = cv2.convert(img, GRAY)}
    \item \relax \texttt{gray = cv2.grayscale(img)}
    \item \relax \texttt{gray = cv2.cvtColor(img, cv2.COLOR\_BGR2GRAY)}
\end{enumerate}
\vspace{2mm}

\noindent\textbf{Q23.} \hfill {\small\color{LPUOrange}\textbf{[OpenCV-based image operations | Medium]}}\\[2pt]
What does the shape attribute \texttt{img.shape} return as \texttt{(480, 640, 3)} indicate about the image?

\begin{enumerate}[label=(\alph*),topsep=2pt,itemsep=1pt,partopsep=0pt,parsep=0pt]
    \item \relax 480 frames, 640 pixels, 3 bits
    \item \relax Height 480, width 640, and 3 color channels
    \item \relax 480 rows, 640 channels, 3 layers
    \item \relax Width 480, height 640, and 3 color channels
\end{enumerate}
\vspace{2mm}

\noindent\textbf{Q24.} \hfill {\small\color{LPUOrange}\textbf{[PyGame-based gaming creation | Medium]}}\\[2pt]
In a typical PyGame program, which function must be called first to start up all imported PyGame modules?

\begin{enumerate}[label=(\alph*),topsep=2pt,itemsep=1pt,partopsep=0pt,parsep=0pt]
    \item \relax \texttt{pygame.start()}
    \item \relax \texttt{pygame.init()}
    \item \relax \texttt{pygame.run()}
    \item \relax \texttt{pygame.begin()}
\end{enumerate}
\vspace{2mm}

\noindent\textbf{Q25.} \hfill {\small\color{LPUOrange}\textbf{[PyGame-based gaming creation | Medium]}}\\[2pt]
Inside a PyGame main loop, which call is responsible for actually updating the full screen with everything drawn?

\begin{enumerate}[label=(\alph*),topsep=2pt,itemsep=1pt,partopsep=0pt,parsep=0pt]
    \item \relax \texttt{pygame.display.init()}
    \item \relax \texttt{pygame.screen.draw()}
    \item \relax \texttt{pygame.display.flip()}
    \item \relax \texttt{pygame.event.get()}
\end{enumerate}
\vspace{2mm}

\noindent\textbf{Q26.} \hfill {\small\color{LPUOrange}\textbf{[PyGame-based gaming creation | Medium]}}\\[2pt]
Why is \texttt{pygame.time.Clock()} combined with \texttt{clock.tick(60)} used in a game loop?

\begin{enumerate}[label=(\alph*),topsep=2pt,itemsep=1pt,partopsep=0pt,parsep=0pt]
    \item \relax To limit the loop to a maximum of 60 frames per second
    \item \relax To pause the game for 60 seconds
    \item \relax To set the background refresh color
    \item \relax To count 60 total game events
\end{enumerate}
\vspace{2mm}

\noindent\textbf{Q27.} \hfill {\small\color{LPUOrange}\textbf{[PyGame-based gaming creation | Medium]}}\\[2pt]
In PyGame, how do you typically detect that the user clicked the window's close button?

\begin{enumerate}[label=(\alph*),topsep=2pt,itemsep=1pt,partopsep=0pt,parsep=0pt]
    \item \relax Call \texttt{pygame.display.quit()} repeatedly
    \item \relax Check for \texttt{event.type == pygame.QUIT} in the event loop
    \item \relax Wait for \texttt{pygame.KEYDOWN}
    \item \relax Check \texttt{pygame.window.closed()}
\end{enumerate}
\vspace{2mm}

\noindent\textbf{Q28.} \hfill {\small\color{LPUOrange}\textbf{[web-scraping-based url access | Medium]}}\\[2pt]
When scraping a webpage in Python, which combination is most commonly used to fetch the HTML and then parse it?

\begin{enumerate}[label=(\alph*),topsep=2pt,itemsep=1pt,partopsep=0pt,parsep=0pt]
    \item \relax \texttt{urllib} to parse, \texttt{requests} to fetch
    \item \relax \texttt{BeautifulSoup} to fetch, \texttt{requests} to parse
    \item \relax \texttt{json} to fetch, \texttt{csv} to parse
    \item \relax \texttt{requests} to fetch, \texttt{BeautifulSoup} to parse
\end{enumerate}
\vspace{2mm}

\noindent\textbf{Q29.} \hfill {\small\color{LPUOrange}\textbf{[web-scraping-based url access | Medium]}}\\[2pt]
After \texttt{response = requests.get(url)}, which attribute do you check to confirm the request succeeded (HTTP 200)?

\begin{enumerate}[label=(\alph*),topsep=2pt,itemsep=1pt,partopsep=0pt,parsep=0pt]
    \item \relax \texttt{response.success}
    \item \relax \texttt{response.ok\_flag}
    \item \relax \texttt{response.result}
    \item \relax \texttt{response.status\_code}
\end{enumerate}
\vspace{2mm}

\noindent\textbf{Q30.} \hfill {\small\color{LPUOrange}\textbf{[web-scraping-based url access | Medium]}}\\[2pt]
Using BeautifulSoup, which method returns a list of all \texttt{<a>} anchor tags on a page?

\begin{enumerate}[label=(\alph*),topsep=2pt,itemsep=1pt,partopsep=0pt,parsep=0pt]
    \item \relax \texttt{soup.find('a')}
    \item \relax \texttt{soup.get('a')}
    \item \relax \texttt{soup.find\_all('a')}
    \item \relax \texttt{soup.select\_one('a')}
\end{enumerate}
\vspace{2mm}

\noindent\textbf{Q31.} \hfill {\small\color{LPUOrange}\textbf{[web-scraping-based url access | Medium]}}\\[2pt]
A website's \texttt{robots.txt} disallows scraping certain paths. What is the correct interpretation for a responsible scraper?

\begin{enumerate}[label=(\alph*),topsep=2pt,itemsep=1pt,partopsep=0pt,parsep=0pt]
    \item \relax It signals which paths crawlers should not access, and should be respected
    \item \relax It automatically blocks all Python requests
    \item \relax It encrypts the page so it cannot be parsed
    \item \relax It is only relevant to search engines and can be ignored
\end{enumerate}
\vspace{2mm}

\noindent\textbf{Q32.} \hfill {\small\color{LPUOrange}\textbf{[data visualization | Medium]}}\\[2pt]
In Matplotlib, which plot type is most appropriate for showing the relationship between two continuous numeric variables?

\begin{enumerate}[label=(\alph*),topsep=2pt,itemsep=1pt,partopsep=0pt,parsep=0pt]
    \item \relax Pie chart
    \item \relax Histogram of one variable
    \item \relax Scatter plot
    \item \relax Bar chart of categories
\end{enumerate}
\vspace{2mm}

\noindent\textbf{Q33.} \hfill {\small\color{LPUOrange}\textbf{[data visualization | Medium]}}\\[2pt]
You want to visualize the distribution and frequency of a single continuous variable such as exam scores. Which chart is best?

\begin{enumerate}[label=(\alph*),topsep=2pt,itemsep=1pt,partopsep=0pt,parsep=0pt]
    \item \relax Histogram
    \item \relax Pie chart
    \item \relax Line chart
    \item \relax Scatter plot
\end{enumerate}
\vspace{2mm}

\noindent\textbf{Q34.} \hfill {\small\color{LPUOrange}\textbf{[data visualization | Medium]}}\\[2pt]
In Matplotlib, what is the purpose of calling \texttt{plt.show()} at the end of a plotting script?

\begin{enumerate}[label=(\alph*),topsep=2pt,itemsep=1pt,partopsep=0pt,parsep=0pt]
    \item \relax It clears all existing plots
    \item \relax It saves the figure to disk
    \item \relax It renders and displays the figure window
    \item \relax It sets the figure resolution
\end{enumerate}
\vspace{2mm}

\noindent\textbf{Q35.} \hfill {\small\color{LPUOrange}\textbf{[data visualization | Medium]}}\\[2pt]
Which library is built on top of Matplotlib and provides high-level functions for attractive statistical graphics?

\begin{enumerate}[label=(\alph*),topsep=2pt,itemsep=1pt,partopsep=0pt,parsep=0pt]
    \item \relax Requests
    \item \relax NumPy
    \item \relax Pillow
    \item \relax Seaborn
\end{enumerate}
\vspace{2mm}

\noindent\textbf{Q36.} \hfill {\small\color{LPUOrange}\textbf{[discussion on advanced toolboxes in python | Medium]}}\\[2pt]
Which Python library is primarily used for numerical computing with efficient multi-dimensional array operations?

\begin{enumerate}[label=(\alph*),topsep=2pt,itemsep=1pt,partopsep=0pt,parsep=0pt]
    \item \relax BeautifulSoup
    \item \relax NumPy
    \item \relax PyGame
    \item \relax Flask
\end{enumerate}
\vspace{2mm}

\noindent\textbf{Q37.} \hfill {\small\color{LPUOrange}\textbf{[discussion on advanced toolboxes in python | Medium]}}\\[2pt]
A data analyst needs to load a CSV file into a labeled, tabular data structure with columns. Which library and object are most appropriate?

\begin{enumerate}[label=(\alph*),topsep=2pt,itemsep=1pt,partopsep=0pt,parsep=0pt]
    \item \relax \texttt{pandas} DataFrame
    \item \relax \texttt{numpy} scalar
    \item \relax \texttt{requests} Session
    \item \relax \texttt{matplotlib} Figure
\end{enumerate}
\vspace{2mm}

\noindent\textbf{Q38.} \hfill {\small\color{LPUOrange}\textbf{[discussion on advanced toolboxes in python | Medium]}}\\[2pt]
Which pair correctly matches a Python toolbox to its main purpose?

\begin{enumerate}[label=(\alph*),topsep=2pt,itemsep=1pt,partopsep=0pt,parsep=0pt]
    \item \relax scikit-learn $\to$ machine learning
    \item \relax Flask $\to$ image processing
    \item \relax TensorFlow $\to$ web scraping
    \item \relax Pandas $\to$ 3D game rendering
\end{enumerate}
\vspace{2mm}

\noindent\textbf{Q39.} \hfill {\small\color{LPUOrange}\textbf{[demonstration of real-world applications | Medium]}}\\[2pt]
A team builds a Python system that detects faces in a live camera feed. Which library is the natural core of this application?

\begin{enumerate}[label=(\alph*),topsep=2pt,itemsep=1pt,partopsep=0pt,parsep=0pt]
    \item \relax OpenCV
    \item \relax Pandas
    \item \relax BeautifulSoup
    \item \relax Flask
\end{enumerate}
\vspace{2mm}

\noindent\textbf{Q40.} \hfill {\small\color{LPUOrange}\textbf{[demonstration of real-world applications | Medium]}}\\[2pt]
A company wants to automatically collect daily product prices from several e-commerce sites into a spreadsheet for analysis. Which workflow best fits?

\begin{enumerate}[label=(\alph*),topsep=2pt,itemsep=1pt,partopsep=0pt,parsep=0pt]
    \item \relax Read the images with OpenCV, then train scikit-learn
    \item \relax Scrape pages with requests/BeautifulSoup, then store in a pandas DataFrame
    \item \relax Render each site with PyGame, then plot with Seaborn
    \item \relax Fetch prices with Matplotlib, then save with NumPy
\end{enumerate}
\vspace{2mm}

\noindent\textbf{Q41.} \hfill {\small\color{LPUOrange}\textbf{[OpenCV-based image operations | Hard]}}\\[2pt]
In OpenCV, you load an image with \texttt{cv2.imread('img.png')} and display it using Matplotlib's \texttt{plt.imshow()}. The colors appear distorted (blues look red). What is the root cause and correct fix?

\begin{enumerate}[label=(\alph*),topsep=2pt,itemsep=1pt,partopsep=0pt,parsep=0pt]
    \item \relax Matplotlib requires normalization; divide the array by 255 before display
    \item \relax OpenCV loads images in BGR order while Matplotlib expects RGB; convert using \texttt{cv2.cvtColor(img, cv2.COLOR\_BGR2RGB)}
    \item \relax The image uses a wrong color depth; convert to grayscale first with \texttt{cv2.COLOR\_BGR2GRAY}
    \item \relax The image is corrupted; reload it with \texttt{cv2.IMREAD\_UNCHANGED}
\end{enumerate}
\vspace{2mm}

\noindent\textbf{Q42.} \hfill {\small\color{LPUOrange}\textbf{[OpenCV-based image operations | Hard]}}\\[2pt]
You apply a 3 $\times$ 3 averaging (box) blur kernel to an image, then apply the same kernel again. Compared to a single 5 $\times$ 5 box blur, the two-pass 3 $\times$ 3 approach:

\begin{enumerate}[label=(\alph*),topsep=2pt,itemsep=1pt,partopsep=0pt,parsep=0pt]
    \item \relax Produces a smoothing effect similar to a 5 $\times$
    \item \relax Produces less smoothing than a single 3 $\times$
    \item \relax Produces exactly the same result as a single 5 $\times$
    \item \relax Cannot be performed because kernels are not associative
\end{enumerate}
\vspace{2mm}

\noindent\textbf{Q43.} \hfill {\small\color{LPUOrange}\textbf{[OpenCV-based image operations | Hard]}}\\[2pt]
When performing edge detection with \texttt{cv2.Canny(img, threshold1, threshold2)}, what best describes the role of the two thresholds in hysteresis?

\begin{enumerate}[label=(\alph*),topsep=2pt,itemsep=1pt,partopsep=0pt,parsep=0pt]
    \item \relax \texttt{threshold1} sets Gaussian blur size and \texttt{threshold2} sets the gradient magnitude
    \item \relax Pixels above \texttt{threshold2} are strong edges; pixels between the thresholds become edges only if connected to strong edges
    \item \relax Both thresholds must be exceeded simultaneously for a pixel to be an edge
    \item \relax Pixels below \texttt{threshold1} are strong edges; those above \texttt{threshold2} are discarded as noise
\end{enumerate}
\vspace{2mm}

\noindent\textbf{Q44.} \hfill {\small\color{LPUOrange}\textbf{[PyGame-based gaming creation | Hard]}}\\[2pt]
In a PyGame loop, a developer moves a sprite by adding a fixed number of pixels each frame without using \texttt{clock.tick()}. On faster hardware the object moves too quickly. What is the correct fix?

\begin{enumerate}[label=(\alph*),topsep=2pt,itemsep=1pt,partopsep=0pt,parsep=0pt]
    \item \relax Call \texttt{pygame.display.flip()} twice per loop iteration
    \item \relax Increase the fixed pixel step to compensate for the frame rate
    \item \relax Multiply movement by delta time (seconds since last frame) so motion is frame-rate independent
    \item \relax Reduce the sprite's image resolution to slow rendering
\end{enumerate}
\vspace{2mm}

\noindent\textbf{Q45.} \hfill {\small\color{LPUOrange}\textbf{[PyGame-based gaming creation | Hard]}}\\[2pt]
A PyGame program's window becomes unresponsive and shows "Not Responding" during gameplay. Which omission most likely causes this?

\begin{enumerate}[label=(\alph*),topsep=2pt,itemsep=1pt,partopsep=0pt,parsep=0pt]
    \item \relax Sprites are drawn before the screen is filled
    \item \relax The display was initialized with too small a resolution
    \item \relax The event queue is never processed with \texttt{pygame.event.get()} or \texttt{pygame.event.pump()}
    \item \relax The frame rate is capped too low with \texttt{clock.tick(30)}
\end{enumerate}
\vspace{2mm}

\noindent\textbf{Q46.} \hfill {\small\color{LPUOrange}\textbf{[PyGame-based gaming creation | Hard]}}\\[2pt]
For pixel-perfect collision between two moving sprites in PyGame, which approach is most accurate but computationally heavier than rectangle collision?

\begin{enumerate}[label=(\alph*),topsep=2pt,itemsep=1pt,partopsep=0pt,parsep=0pt]
    \item \relax Checking Euclidean distance between sprite centers
    \item \relax Using \texttt{pygame.sprite.spritecollide} with \texttt{collided=None}
    \item \relax Mask-based collision using \texttt{pygame.mask.from\_surface} and \texttt{mask.overlap}
    \item \relax Comparing \texttt{rect.colliderect} on both sprites
\end{enumerate}
\vspace{2mm}

\noindent\textbf{Q47.} \hfill {\small\color{LPUOrange}\textbf{[web-scraping-based url access | Hard]}}\\[2pt]
You scrape a page with \texttt{requests.get(url)} and BeautifulSoup, but the target data is missing from the returned HTML even though it appears in the browser. What is the most probable reason?

\begin{enumerate}[label=(\alph*),topsep=2pt,itemsep=1pt,partopsep=0pt,parsep=0pt]
    \item \relax The content is rendered client-side by JavaScript after page load, which \texttt{requests} does not execute
    \item \relax BeautifulSoup cannot parse \texttt{<div>} tags without the \texttt{lxml} parser
    \item \relax \texttt{requests} strips all dynamic tags from the response by default
    \item \relax The server returned a 200 status so the data must be present in a hidden attribute
\end{enumerate}
\vspace{2mm}

\noindent\textbf{Q48.} \hfill {\small\color{LPUOrange}\textbf{[web-scraping-based url access | Hard]}}\\[2pt]
Which statement about \texttt{robots.txt} and web scraping is most accurate?

\begin{enumerate}[label=(\alph*),topsep=2pt,itemsep=1pt,partopsep=0pt,parsep=0pt]
    \item \relax It is an advisory standard that requests crawlers avoid certain paths, but it does not technically enforce access control
    \item \relax It is a firewall rule that blocks all unauthorized scraping requests
    \item \relax It automatically rate-limits any client that ignores it
    \item \relax It encrypts URLs so scrapers cannot read the site structure
\end{enumerate}
\vspace{2mm}

\noindent\textbf{Q49.} \hfill {\small\color{LPUOrange}\textbf{[web-scraping-based url access | Hard]}}\\[2pt]
A scraper repeatedly hits a site and starts receiving HTTP 429 responses. Which combination best addresses this while remaining respectful?

\begin{enumerate}[label=(\alph*),topsep=2pt,itemsep=1pt,partopsep=0pt,parsep=0pt]
    \item \relax Switch every request to a new random User-Agent and ignore the status code
    \item \relax Send requests faster to finish before the block takes effect
    \item \relax Disable SSL verification to bypass the rate limiter
    \item \relax Add delays between requests and implement exponential backoff honoring the \texttt{Retry-After} header
\end{enumerate}
\vspace{2mm}

\noindent\textbf{Q50.} \hfill {\small\color{LPUOrange}\textbf{[data visualization | Hard]}}\\[2pt]
You plot data spanning 1 to 1 0 6 on a linear y-axis and small values become invisible near zero. Which change best reveals structure across all magnitudes?

\begin{enumerate}[label=(\alph*),topsep=2pt,itemsep=1pt,partopsep=0pt,parsep=0pt]
    \item \relax Increase the figure DPI to 300
    \item \relax Use a logarithmic y-axis (e.g., \texttt{plt.yscale('log')})
    \item \relax Add more gridlines to the linear axis
    \item \relax Switch the plot color to a brighter palette
\end{enumerate}
\vspace{2mm}

\noindent\textbf{Q51.} \hfill {\small\color{LPUOrange}\textbf{[data visualization | Hard]}}\\[2pt]
When designing a chart for viewers who may be color-blind, which practice is most effective?

\begin{enumerate}[label=(\alph*),topsep=2pt,itemsep=1pt,partopsep=0pt,parsep=0pt]
    \item \relax Combine a color-blind-safe palette with redundant encodings like markers, line styles, or direct labels
    \item \relax Rely on a rainbow (jet) colormap for continuous data
    \item \relax Remove the legend so viewers focus on colors
    \item \relax Use only red and green to maximize contrast
\end{enumerate}
\vspace{2mm}

\noindent\textbf{Q52.} \hfill {\small\color{LPUOrange}\textbf{[data visualization | Hard]}}\\[2pt]
A bar chart of average sales per region hides the fact that one region has a bimodal distribution. Which visualization better exposes this within-group structure?

\begin{enumerate}[label=(\alph*),topsep=2pt,itemsep=1pt,partopsep=0pt,parsep=0pt]
    \item \relax A violin plot or box plot showing the distribution shape per region
    \item \relax A pie chart of the same averages
    \item \relax A line chart connecting the regional means
    \item \relax A single stacked bar of all regions combined
\end{enumerate}
\vspace{2mm}

\noindent\textbf{Q53.} \hfill {\small\color{LPUOrange}\textbf{[data visualization | Hard]}}\\[2pt]
In Matplotlib, what is the key difference between the stateful \texttt{pyplot} interface and the object-oriented (Axes) interface for complex multi-panel figures?

\begin{enumerate}[label=(\alph*),topsep=2pt,itemsep=1pt,partopsep=0pt,parsep=0pt]
    \item \relax The \texttt{pyplot} interface renders faster because it caches figures
    \item \relax Only the \texttt{pyplot} interface supports saving figures to file
    \item \relax The OO interface gives explicit control over each \texttt{Axes} object, avoiding ambiguity about which subplot commands affect
    \item \relax The OO interface cannot create subplots without \texttt{plt.figure()}
\end{enumerate}
\vspace{2mm}

\noindent\textbf{Q54.} \hfill {\small\color{LPUOrange}\textbf{[discussion on advanced toolboxes in python | Hard]}}\\[2pt]
You perform an element-wise operation on two large arrays. Why is NumPy vectorization typically much faster than a Python \texttt{for} loop?

\begin{enumerate}[label=(\alph*),topsep=2pt,itemsep=1pt,partopsep=0pt,parsep=0pt]
    \item \relax Python \texttt{for} loops cannot operate on numeric data
    \item \relax NumPy stores data on the GPU by default
    \item \relax NumPy executes the operation in compiled C code over contiguous memory, avoiding per-element Python interpreter overhead
    \item \relax NumPy loops are automatically parallelized across all CPU cores by default
\end{enumerate}
\vspace{2mm}

\noindent\textbf{Q55.} \hfill {\small\color{LPUOrange}\textbf{[discussion on advanced toolboxes in python | Hard]}}\\[2pt]
A Pandas operation \texttt{df[df['a'] > 0]['b'] = 5} produces a \texttt{SettingWithCopyWarning} and the change may not persist. What is the correct fix?

\begin{enumerate}[label=(\alph*),topsep=2pt,itemsep=1pt,partopsep=0pt,parsep=0pt]
    \item \relax Disable the warning with \texttt{pd.options.mode.chained\_assignment = None}
    \item \relax Use \texttt{.loc} for combined selection and assignment: \texttt{df.loc[df['a'] > 0, 'b'] = 5}
    \item \relax Convert the column to a list, modify it, and reassign
    \item \relax Wrap the expression in \texttt{df.copy()} before assigning
\end{enumerate}
\vspace{2mm}

\noindent\textbf{Q56.} \hfill {\small\color{LPUOrange}\textbf{[discussion on advanced toolboxes in python | Hard]}}\\[2pt]
When training a scikit-learn model, why should scaling (e.g., \texttt{StandardScaler}) be fit only on the training data and not on the full dataset before splitting?

\begin{enumerate}[label=(\alph*),topsep=2pt,itemsep=1pt,partopsep=0pt,parsep=0pt]
    \item \relax Fitting on all data leaks test-set statistics into training, producing overly optimistic evaluation
    \item \relax The scaler cannot process both sets simultaneously due to memory limits
    \item \relax Test data must never be numerically transformed
    \item \relax Scaling the test set is impossible without the training labels
\end{enumerate}
\vspace{2mm}

\noindent\textbf{Q57.} \hfill {\small\color{LPUOrange}\textbf{[storytelling | Hard]}}\\[2pt]
In data storytelling, a stakeholder audience needs to grasp one key takeaway quickly. Which approach best supports this goal?

\begin{enumerate}[label=(\alph*),topsep=2pt,itemsep=1pt,partopsep=0pt,parsep=0pt]
    \item \relax Remove all text so the numbers speak for themselves
    \item \relax Present every chart produced during analysis in chronological order
    \item \relax Maximize data density by overlaying all metrics on one axis
    \item \relax Lead with a clear headline insight and use annotations to direct attention to the supporting evidence
\end{enumerate}
\vspace{2mm}

\noindent\textbf{Q58.} \hfill {\small\color{LPUOrange}\textbf{[storytelling | Hard]}}\\[2pt]
A chart truncates its y-axis to start at 95 instead of 0 , making a 2\% change look dramatic. In the context of ethical storytelling, this is best described as:

\begin{enumerate}[label=(\alph*),topsep=2pt,itemsep=1pt,partopsep=0pt,parsep=0pt]
    \item \relax A required convention for all bar charts
    \item \relax A best practice that improves resolution and is always appropriate
    \item \relax A misleading design choice that exaggerates differences and should be disclosed or avoided
    \item \relax An error that Matplotlib prevents automatically
\end{enumerate}
\vspace{2mm}

\noindent\textbf{Q59.} \hfill {\small\color{LPUOrange}\textbf{[demonstration of real-world applications | Hard]}}\\[2pt]
You must process a live video stream, detect motion, and log events to a dashboard. Which combination of Python tools maps most appropriately to each subtask?

\begin{enumerate}[label=(\alph*),topsep=2pt,itemsep=1pt,partopsep=0pt,parsep=0pt]
    \item \relax Requests for capture, scikit-learn for logging, and PyGame for the dashboard
    \item \relax OpenCV for frame capture/motion detection, Pandas for event logging, and Matplotlib/Plotly for the dashboard
    \item \relax Matplotlib for capture, OpenCV for logging, and Requests for the dashboard
    \item \relax PyGame for capture, BeautifulSoup for motion detection, and NumPy for the dashboard
\end{enumerate}
\vspace{2mm}

\noindent\textbf{Q60.} \hfill {\small\color{LPUOrange}\textbf{[demonstration of real-world applications | Hard]}}\\[2pt]
A pipeline scrapes prices hourly, stores them, and predicts trends. To keep it maintainable and avoid re-scraping unchanged pages, which design choice is most sound?

\begin{enumerate}[label=(\alph*),topsep=2pt,itemsep=1pt,partopsep=0pt,parsep=0pt]
    \item \relax Combine scraping, storage, and modeling into one function for simplicity
    \item \relax Store all data only in memory to speed up predictions
    \item \relax Cache responses with conditional requests (ETag/If-Modified-Since) and separate scraping, storage, and modeling into modular stages
    \item \relax Scrape every page fully each run regardless of changes to ensure freshness
\end{enumerate}
\vspace{2mm}


\newpage
\section{Section B: Comprehensive Subjective Problem with Detailed Analytical Rubric}

\begin{subjectiveq}{6}{1}
Design an applied engineering pipeline integrating OpenCV image processing and web API data retrieval:
\begin{enumerate}[label=(\alph*)]
    \item Write an OpenCV script that reads an image, converts BGR to Grayscale, applies Gaussian Blur with $5 \times 5$ kernel to eliminate high-frequency noise, detects edges using Canny edge detector ($T_{\text{lower}}=50, T_{\text{upper}}=150$), and finds contours.
    \item Implement a resilient Python web scraper using \texttt{requests} and \texttt{BeautifulSoup} that extracts real-time product prices from an e-commerce page, handles HTTP 429 rate limiting with exponential backoff, and saves structured records into JSON format.
\end{enumerate}
\end{subjectiveq}

\begin{solbox}
\textbf{OpenCV \& Web Scraping Implementations:}
\begin{pycode}
import cv2\\
import time, requests, json\\
from bs4 import BeautifulSoup\\
~\\
def process\_image\_contours(image\_path):\\
\quad img = cv2.imread(image\_path)\\
\quad if img is None:\\
\quad \quad raise FileNotFoundError(f"Cannot read image at \{image\_path\}")\\
\quad gray = cv2.cvtColor(img, cv2.COLOR\_BGR2GRAY)\\
\quad blurred = cv2.GaussianBlur(gray, (5, 5), 1.4)\\
\quad edges = cv2.Canny(blurred, 50, 150)\\
\quad contours, \_ = cv2.findContours(edges, cv2.RETR\_EXTERNAL, cv2.CHAIN\_APPROX\_SIMPLE)\\
\quad return len(contours), edges\\
~\\
def fetch\_with\_backoff(url, max\_retries=3):\\
\quad headers = \{"User-Agent": "AcademicResearcher/1.0"\}\\
\quad for attempt in range(max\_retries):\\
\quad \quad resp = requests.get(url, headers=headers)\\
\quad \quad if resp.status\_code == 200:\\
\quad \quad \quad soup = BeautifulSoup(resp.text, 'html.parser')\\
\quad \quad \quad return soup\\
\quad \quad elif resp.status\_code == 429:\\
\quad \quad \quad time.sleep(2 ** attempt)\\
\quad raise RuntimeError("Max retries exceeded or forbidden")
\end{pycode}
\textbf{Engineering Rubric:} Highlight proper channel ordering (BGR vs RGB), threshold selection rationale for Canny suppression, header spoofing ethics, and robots.txt compliance.
\end{solbox}

\begin{examinertip}
\begin{itemize}[leftmargin=1.5em, itemsep=2pt]
    \item \textbf{Algorithmic Formulation (3 Marks):} Correct problem decomposition, function signatures, and parameter documentation.
    \item \textbf{Syntactic Rigor \& Implementation (4 Marks):} Valid Python syntax, proper error handling, defensive input validation, and avoidance of anti-patterns.
    \item \textbf{Complexity \& Edge Analysis (3 Marks):} Accurate time and space asymptotic analysis $\mathcal{O}(n)$ and testing edge conditions (null inputs, zero division, empty collections).
\end{itemize}
\end{examinertip}
